% Consolidated 05-rigid-body from 6915d31494b4647dd12b50a59cddbc000c3484fc; original TeX fragments preserved.

% === BEGIN 05-rigid-body-234-240.tex blob 0d90c14765244a7eee61dc21dc67e9e6f244c4de ===
% Irodov user PDF p.46--47, original printed pp.37--38.
% Opening of section 1.5: rigid-body statics and inertia tensors.
\exproblem{234}{37}
一根质量为 $m=1.0\,\mathrm{kg}$ 的均匀细杆 $AB$
在同一平面内受两个反平行力作用而只作平动，
其加速度为 $a_C=2.0\,\mathrm{m/s^2}$（图1.52）。
两个施力点沿杆方向相距
$d=20\,\mathrm{cm}$，已知其中一个力
$F_2=5.0\,\mathrm N$，求杆长 $\ell$。
\exsolution
按原图，上端 $A$ 的力 $F_1$ 向左，
另一力 $F_2$ 向右且施力点比 $A$ 低 $d$。
向右取正，纯平动要求
$F_2-F_1=ma_C$，故
\[
F_1=5-1(2)=3.0\,\mathrm N.
\]
均匀细杆质心在杆中点。
无角加速度意味着关于质心的合力矩为零：
\[
F_1\frac{\ell}{2}
=F_2\left(\frac{\ell}{2}-d\right).
\]
移项得
\[
\boxed{\ell=\frac{2F_2d}{F_2-F_1}
=\frac{2(5)(0.20)}{2.0}=1.00\,\mathrm m}.
\]
\excheck 当所有力都沿水平方向时，
平动加速度由合力决定；无转动则必须另外满足
关于质心合力矩为零。

\exproblem{235}{37}
在 $xy$ 平面中，力
$\bm F=A\bm i+B\bm j$ 作用于位置
$\bm r=a\bm i+b\bm j$。
求该力关于坐标原点 $O$ 的力矩矢量及力臂长度。
\exsolution
力矩由叉积定义：
\[
\bm N_O=\bm r\times\bm F
=(a\bm i+b\bm j)\times(A\bm i+B\bm j)
=(aB-bA)\bm k.
\]
力臂是从原点到作用线的垂直距离，
利用 $|\bm N|=|\bm F|\ell$：
\[
\boxed{\bm N_O=(aB-bA)\bm k},\qquad
\boxed{\ell=\frac{|aB-bA|}{\sqrt{A^2+B^2}}}.
\]
\excheck 若 $aB=bA$，力作用线穿过原点，
力矩和力臂均为零；若 $\bm F=\bm0$，
则作用线未定义而力矩仍为零。

\exproblem{236}{37}
在平面内两个力分别为
$\bm F_1=A\bm j$、$\bm F_2=B\bm i$，
施力点的位置矢量分别为
$\bm r_1=a\bm i$ 和 $\bm r_2=b\bm j$。
求合力相对于原点的力臂。
\exsolution
合力矢量
$\bm F=\bm F_1+\bm F_2=B\bm i+A\bm j$。
两个力矩分别为
$\bm N_1=aA\bm k$ 与
$\bm N_2=-bB\bm k$，故总力矩
\[
\bm N=(aA-bB)\bm k.
\]
因此合力大小为 $\sqrt{A^2+B^2}$ 时，
其等效作用线到原点的距离
\[
\boxed{\ell=\frac{|aA-bB|}{\sqrt{A^2+B^2}}}.
\]
\excheck 力臂必须由\emph{合力的总力矩}计算；
不能把两力各自的力臂直接相加。

\exproblem{237}{37}
如图1.53，正方形薄板顶点依次为
$A$（左）、$B$（上）、$C$（右）、$D$（下）。
在 $A$ 点施竖直向上的力 $F$，
在 $B$ 点施水平向右的力 $F$，
在 $D$ 点沿右下方向施大小为 $F\sqrt2$ 的力。
求合力的大小、方向；
若作用点选择在边 $BC$ 上，求它的位置。
\exsolution
取正方形中心为原点，四顶点分别为
$A=(-d,0),B=(0,d),C=(d,0),D=(0,-d)$。
三个力的坐标分量分别为
$(0,F),(F,0),(F,-F)$。故
\[
\bm F_{\rm res}=(2F,0),
\quad\boxed{|\bm F_{\rm res}|=2F},
\]
方向水平向右，即平行于对角线 $AC$。
总力矩（取垂直板面分量）为
\[
N_z=(-d)F-dF+dF=-dF.
\]
设等效合力作用点为 $(x,y)$，
则 $N_z=xF_y-yF_x=-2Fy$，
所以 $y=d/2$。
边 $BC$ 上从 $(0,d)$ 到 $(d,0)$
满足 $x+y=d$，于是 $x=d/2$。
\[
\boxed{\text{合力作用于边 }BC\text{ 的中点。}}
\]
\excheck 该解保留了总力和总力矩，因此与原三力体系
对薄板的平动、转动作用等效。

\exproblem{238}{37--38}
求：（a）长为 $\ell$、质量 $m$ 的均匀细杆
对过一端且垂直于杆的轴的转动惯量；
（b）长宽为 $a,b$、质量为 $m$ 的均匀矩形薄板
对过一个顶点且垂直于薄板的轴的转动惯量。
\exsolution
（a）细杆线密度为 $m/\ell$，
从转轴端量距离 $x\in[0,\ell]$，
\[
I_a=\int_0^\ell x^2\frac m\ell\,\mathrm dx
=\boxed{\frac13m\ell^2}.
\]
（b）选顶点为原点、两邻边沿坐标轴，
薄板面密度为 $m/(ab)$，轴距平方为 $x^2+y^2$：
\[
I_b=\frac{m}{ab}\int_0^a\int_0^b
(x^2+y^2)\,\mathrm dy\,\mathrm dx
=\boxed{\frac m3(a^2+b^2)}.
\]
\excheck 也可由平行轴定理验证，矩形板的质心轴惯量
$m(a^2+b^2)/12$，顶点距质心平方为 $(a^2+b^2)/4$，
相加得到同一结果。

\exproblem{239}{38}
求：（a）均匀铜圆盘，厚度
$b=2.0\,\mathrm{mm}$，半径
$R=100\,\mathrm{mm}$，绕轴线的转动惯量；
（b）质量 $m$、底面半径 $R$ 的均匀实心圆锥
绕其轴线的转动惯量。
\exsolution
（a）铜密度记为 $\rho$，
圆盘质量为 $M=\rho\pi R^2b$。
由均匀实心圆盘绕轴惯量
$I=\tfrac12MR^2$，得
\[
\boxed{I_a=\frac12\pi\rho bR^4}.
\]
取铜密度 $\rho\simeq8.9\times10^3\,\mathrm{kg/m^3}$，
数值得
$\boxed{I_a\simeq2.80\times10^{-3}\,\mathrm{kg\,m^2}}$。
（b）设圆锥高为 $H$，从锥顶起截面半径
$r(z)=Rz/H$。在高度 $z$ 处取厚度
$\mathrm dz$ 的薄圆盘，
$\mathrm dm=\rho\pi r(z)^2\,\mathrm dz$，
微元绕轴惯量 $\mathrm dI=\tfrac12r(z)^2\,\mathrm dm$。
积分：
\[
I=\frac12\rho\pi\frac{R^4}{H^4}
\int_0^H z^4\,\mathrm dz
=\frac{\rho\pi R^4H}{10}.
\]
又 $m=\rho\pi R^2H/3$，故
\[
\boxed{I_b=\frac{3}{10}mR^2}.
\]
\excheck $I_a$ 的原书单位写法 $2.8\,\mathrm{g\,m^2}$
与 $2.8\times10^{-3}\,\mathrm{kg\,m^2}$ 相同；
此处统一用 SI 单位。

\exproblem{240}{38}
证明：任意形状的薄板通过同一点的三条互相垂直轴中，
若前两条轴在板面内、第三条轴垂直板面，
则转动惯量满足 $I_1+I_2=I_3$。
据此求半径 $R$、质量 $m$ 的均匀薄圆盘
绕任一直径的转动惯量。
\exsolution
选薄板面为 $xy$ 平面，三条轴分别为 $x,y,z$，
薄板任一质元到相应轴的距离平方分别为
$y^2,x^2,x^2+y^2$。按转动惯量定义：
\[
I_x=\int y^2\,\mathrm dm,\quad
I_y=\int x^2\,\mathrm dm,\quad
I_z=\int(x^2+y^2)\,\mathrm dm.
\]
积分线性性直接给出
\[
\boxed{I_x+I_y=I_z}.
\]
均匀圆盘的对称性意味着过圆心的两条正交直径轴
转动惯量相等：$I_x=I_y$。
由于 $I_z=\tfrac12mR^2$，
\[
\boxed{I_{\rm diameter}=\frac12I_z=\frac14mR^2}.
\]
\excheck 垂直轴定理只适用于质量分布被认为在同一平面内的
薄板；对厚度不能忽略的三维刚体不能直接套用。
% === END 05-rigid-body-234-240.tex ===

% === BEGIN 05-rigid-body-241-250.tex blob 560d2a036b0ac004307d023cc957bc4f9fa6045e ===
% I. E. Irodov, Problems in General Physics, user-provided Arihant PDF.
% Problems 241--250; original printed pages 38--39; PDF pages 47--48.
% First draft of Chinese translation with independently worked equations.

\exproblem{241}{38}
一块原本均匀、半径为 $R=20\,\mathrm{cm}$ 的圆盘，
按原书图1.54挖去一个直径为 $R$ 的圆孔：
孔与原盘外缘内切，孔边缘经过原圆心 $O$。
剩余阴影部分质量为 $m=7.3\,\mathrm{kg}$。
求剩余部分对通过其\emph{质心且垂直盘面}的轴的转动惯量。
\exsolution
设原盘面密度为 $\sigma$，孔半径 $R/2$，
孔心在 $x=-R/2$。原盘与被挖去圆片的面积比为 $4:1$，
故其质量分别为 $M=4m/3$、$m_h=m/3$。
取原圆心为 $x=0$，剩余部分质心位置满足
\[
mx_C=M(0)-m_h(-R/2)
\quad\Longrightarrow\quad
\boxed{x_C=R/6}.
\]
对原圆心 $O$ 的惯量可用“完整盘减去孔”计算。
平行轴定理给出被挖圆片对 $O$ 的惯量为
\[
I_{h,O}=\tfrac12m_h(R/2)^2+m_h(R/2)^2
=\tfrac38m_h R^2.
\]
因此
\[
I_O=\tfrac12MR^2-I_{h,O}
=\left(\frac23-\frac18\right)mR^2
=\frac{13}{24}mR^2.
\]
再向剩余部分的质心轴平移：
\[
\boxed{I_C=I_O-mx_C^2
=\left(\frac{13}{24}-\frac1{36}\right)mR^2
=\frac{37}{72}mR^2
\simeq0.150\,\mathrm{kg\,m^2}.}
\]
\excheck 这里的 $m$ 是\emph{挖孔后}剩余部分质量，
不是完整圆盘质量。所得数值与底本书末答案相符。

\exproblem{242}{38}
利用均匀实心球的转动惯量公式，
求质量为 $m$、半径为 $R$ 的均匀薄球壳
关于过球心直径轴的转动惯量。
\exsolution
均匀实心球的质量、转动惯量分别为
$M(R)=\frac43\pi\rho R^3$、
$I(R)=\frac25M(R)R^2=\frac8{15}\pi\rho R^5$。
设把半径从 $R$ 增加到 $R+\mathrm dR$，
新增薄层质量及转动惯量为
\[
\mathrm dm=\frac{\mathrm dM}{\mathrm dR}\mathrm dR
=4\pi\rho R^2\,\mathrm dR,\qquad
\mathrm dI=\frac{\mathrm dI}{\mathrm dR}\mathrm dR
=\frac83\pi\rho R^4\,\mathrm dR.
\]
二者比值
$\mathrm dI/\mathrm dm=\frac23R^2$，
令薄层总质量为 $m$，得到
\[
\boxed{I_{\mathrm{shell}}=\frac23mR^2}.
\]
\excheck 球壳的 $2/3$ 不等于均匀实心球的 $2/5$；
质量都集中于外半径附近时惯量更大。

\exproblem{243}{38}
质量为 $M$、半径为 $R$ 的均匀实心圆柱
可绕固定水平轴无摩擦转动。
圆柱上缠绕轻绳，绳自由端悬挂质量 $m$ 的重物（图1.55）。
系统从静止释放，求：(a) 圆柱角速度 $\omega(t)$；
(b) 包括重物在内的系统总动能 $K(t)$。
\exsolution
以重物向下为正，绳无滑动，故
$a=R\beta$。记绳张力为 $T$，圆柱惯量
$I=\frac12MR^2$。牛顿方程和转动方程为
\[
mg-T=ma,\qquad TR=I\beta=\frac12MRa.
\]
于是 $T=\frac12Ma$，消元得
\[
a=\frac{mg}{m+M/2}=\frac{g}{1+M/(2m)}.
\]
由初始静止和 $\beta=a/R$，
\[
\boxed{\omega(t)=\frac{gt}{R[1+M/(2m)]}}.
\]
重物速度 $v=at$，所以总动能
\[
K=\frac12mv^2+\frac12I(v/R)^2
=\frac12(m+M/2)a^2t^2.
\]
化简后
\[
\boxed{K(t)=\frac{mg^2t^2}{2[1+M/(2m)]}}.
\]
\excheck 两式与原书简答相符；
$m\to\infty$（相对圆柱）时 $a\to g$，
$M\to\infty$ 时 $a\to0$。

\exproblem{244}{38}
麦克斯韦轮的轮轴半径为 $r$，
轮和轮轴总质量为 $m$，绕对称轴转动惯量为 $I$。
两根轻绳对称缠绕于轮轴，
其上端固定在可竖直升降的横梁上。
横梁以恰当的竖直加速度向上运动，使轮的质心高度不变，
同时绳逐渐放开。求：(a) 每根绳的张力；
(b) 横梁向上的加速度。
\exsolution
轮质心始终不作竖直加速，故竖直力平衡
\[
2T-mg=0,\qquad
\boxed{T=\frac{mg}{2}}.
\]
两根绳对轮轴施加同向力矩
$\tau=2Tr= mgr$。
轮绕自身轴的角加速度
$\beta=\tau/I=mgr/I$。
绳与轮轴间不滑动，轮心静止时横梁向上的
线加速度大小等于轮轴表面放绳切向加速度：
\[
\boxed{a_{\rm beam}=r\beta=\frac{mgr^2}{I}}.
\]
\excheck 横梁提供能量使轮加速自转，而轮心高度不变。
该结论的 $I$ 必须包括整个轮与轮轴。

\exproblem{245}{38}
一根质量为 $m$、长度为 $\ell$ 的均匀水平细杆，
可绕通过一端 $A$ 的固定竖轴无摩擦转动。
原来静止，另一端 $B$ 从此时起受到水平恒力 $F$，
力的方向\emph{始终垂直于杆的原始方向}，而非始终垂直于杆。
求角速度随杆转过的角度 $\varphi$ 的关系。
\exsolution
杆关于端点轴的转动惯量
$I_A=m\ell^2/3$。由于力矢量在实验室坐标中
方向固定，杆端沿该方向的位移为
$\Delta y_B=\ell\sin\varphi$，因此外力做功
$W=F\ell\sin\varphi$。
若杆从 $\varphi=0$ 静止开始且没有损耗，
动能定理给出
\[
\tfrac12I_A\omega^2=F\ell\sin\varphi.
\]
故
\[
\boxed{\omega(\varphi)=
\sqrt{\frac{6F\sin\varphi}{m\ell}}}
\]
在从原方向开始加速的 $0\le\varphi\le\pi$
运动段取正向角速度。
\excheck 若错误认为力始终垂直当前杆，
会得到 $\omega^2\propto\varphi$，
与题面“perpendicular to the original position”不符。
底本 PDF 文本抽取可能漏识别答案中的平方根。

\exproblem{246}{39}
如原书图1.56，两个悬挂物体质量分别为
$m_1,m_2$，连接轻绳跨过半径 $R$、
质量 $m$ 的均匀实心圆柱滑轮。
绳与滑轮间无相对滑动，轴无摩擦。
求：(a) 滑轮角加速度；
(b) 两侧绳张力比 $T_1/T_2$。
\exsolution
先取 $m_2>m_1$、物体2向下加速度为 $a$，
绳长约束使物体1向上同加速度。
记滑轮惯量 $I=\frac12mR^2$。列方程
\[
T_1-m_1g=m_1a,\qquad
m_2g-T_2=m_2a,\qquad
(T_2-T_1)R=I\frac aR.
\]
代入张力并化简：
\[
a=\frac{(m_2-m_1)g}{m_1+m_2+m/2},
\quad
\boxed{|\beta|=\frac{|m_2-m_1|g}
{R(m_1+m_2+m/2)}}.
\]
此外
\[
T_1=m_1(g+a)
=\frac{m_1(m+4m_2)g}{2(m_1+m_2+m/2)},
\quad
T_2=m_2(g-a)
=\frac{m_2(m+4m_1)g}{2(m_1+m_2+m/2)}.
\]
故
\[
\boxed{\frac{T_1}{T_2}
=\frac{m_1(m+4m_2)}{m_2(m+4m_1)}}.
\]
\excheck 当 $m_1=m_2$ 时 $\beta=0$ 且张力比为1；
当滑轮质量趋零时亦恢复轻滑轮的相等张力。

\exproblem{247}{39}
质量为 $m_1$ 的物体在摩擦系数为 $k$ 的
水平面上，轻绳跨过质量为 $m$、
可视为均匀圆盘的有质量滑轮，另一端悬挂 $m_2$（图1.57）。
系统从静止释放，且已知 $m_2$ 开始下降。
忽略绳质量与轴摩擦，求最初 $t$ 秒摩擦力对 $m_1$ 做的功。
\exsolution
设 $m_2$ 向下、$m_1$ 向右的加速度均为 $a$，
滑轮半径为 $R$，惯量 $I=\frac12mR^2$。
将三部分的动力学方程相加，绳张力在整体中抵消：
\[
m_2g-km_1g=(m_1+m_2+I/R^2)a.
\]
故
\[
a=\frac{(m_2-km_1)g}{m_1+m_2+m/2}>0.
\]
前 $t$ 秒物体1位移 $s=\frac12at^2$，
滑动摩擦力大小 $km_1g$ 且反向，
\[
W_f=-km_1gs
=\boxed{-\frac{k m_1(m_2-km_1)g^2t^2}
{m+2(m_1+m_2)}}.
\]
\excheck 题面声明 $m_2$ 开始下降，等价于
$ m_2>km_1$ 的滑动模型阈值，故最终功为负。

\exproblem{248}{39}
半径 $R$ 的均匀实心圆柱绕自身轴以初角速度
$\omega_0$ 旋转，随后放入两互相垂直的粗糙墙形成的角落，
如原图1.58。
两接触面摩擦系数同为 $k$，圆柱中心始终保持静止。
求圆柱停止前转过多少圈 $n$。
\exsolution
取原图所示旋向，底部摩擦力水平向左，
左壁摩擦力竖直向上。记底部支持力 $N_b$、
左壁支持力 $N_w$。
质心无平动，水平与竖直方向的力平衡给出
\[
N_w=kN_b,\qquad N_b+kN_w=mg.
\]
所以
$N_b=mg/(1+k^2)$、
$N_w=kmg/(1+k^2)$。
两处滑动摩擦对中心产生同向阻转力矩：
\[
\tau_f=kR(N_b+N_w)
=\frac{k(1+k)mgR}{1+k^2}.
\]
实心圆柱惯量 $I=\frac12mR^2$，
角减速度大小为
$\beta_f=\tau_f/I=
2k(1+k)g/[(1+k^2)R]$。
用 $\omega_0^2=2\beta_f\Delta\varphi$，
并以 $n=\Delta\varphi/(2\pi)$ 换算圈数：
\[
\boxed{n=\frac{(1+k^2)\omega_0^2R}
{8\pi k(1+k)g}}.
\]
\excheck 此结果依赖图示旋向；反向旋转时两面摩擦
竖直、水平分量的组合改变，应重新列平衡方程，
不能不检查受力方向就套用此式。

\exproblem{249}{39}
半径为 $R$ 的均匀圆盘绕竖直中心轴以角速度
$\omega_0$ 旋转，然后轻放在水平粗糙面上。
摩擦系数为 $k$，盘对地的压强近似均匀。
求圆盘继续旋转多久才停止。
\exsolution
圆盘总支持力 $mg$，均匀压强
$p=mg/(\pi R^2)$。半径为 $r$、
宽度为 $\mathrm dr$ 的圆环与地面的接触面积为
$2\pi r\,\mathrm dr$，其摩擦力为
$k p\,2\pi r\,\mathrm dr$，对中心的阻转力矩微元
\[
\mathrm d\tau=r\cdot k p\,2\pi r\,\mathrm dr
=2\pi kp r^2\,\mathrm dr.
\]
积分得 $\tau=\frac23kmgR$。
圆盘惯量 $I=\frac12mR^2$，
角减速度恒为
$\beta=\tau/I=4kg/(3R)$。
于是从角速度 $\omega_0$ 减到零所需时间
\[
\boxed{t_{\rm stop}=\frac{\omega_0}{\beta}
=\frac{3\omega_0R}{4kg}}.
\]
\excheck 若误把所有接触摩擦集中作用在边缘 $r=R$，
会把阻转力矩错误估成 $kmgR$，结果偏小。

\exproblem{250}{39}
一飞轮初始角速度为 $\omega_0$，
受到大小与瞬时角速度平方根成正比的阻转力矩：
$|\tau|=c\sqrt{\omega}$，其中 $c>0$。
求整个减速停止阶段按时间平均的角速度。
\exsolution
设飞轮转动惯量为 $I$，
转动方程
$I\dot\omega=-c\sqrt{\omega}$。
分离变量积分：
\[
2\sqrt{\omega(t)}=
2\sqrt{\omega_0}-\frac cI t.
\]
停止时间为 $T=2I\sqrt{\omega_0}/c$，
因此在 $0\le t\le T$，
\[
\omega(t)=\omega_0\left(1-\frac tT\right)^2.
\]
求时间平均：
\[
\boxed{\bar\omega=\frac1T\int_0^T
\omega_0(1-t/T)^2\,\mathrm dt
=\omega_0\int_0^1(1-u)^2\,\mathrm du
=\frac{\omega_0}{3}}.
\]
\excheck 与角速度衰减成线性时的 $\omega_0/2$ 不同；
阻转系数 $c$ 只决定停转时间，不影响上述无量纲平均值。
% === END 05-rigid-body-241-250.tex ===

% === BEGIN 05-rigid-body-251-260.tex blob 915c6ad2d7d1b4715e20a7daffd24b9a666ce101 ===
% Irodov, user-provided PDF, printed pp. 39--41; PDF pp. 48--50.
% Original diagrams 1.59--1.66 were inspected for pulley, spool and rolling constraints.

\exproblem{251}{39}
半径为 $R$、质量为 $M$ 的均匀实心圆柱可绕固定的水平轴
$O$ 无摩擦转动。质量为 $m$、总长为 $\ell$ 的均匀轻柔细绳
沿圆柱表面缠绕成一层（图1.59），自由下垂部分长为 $x$。
题目规定缠绕部分的质心可视为在圆柱轴上。
求圆柱的角加速度 $\beta(x)$。
\exsolution
记绳线密度为 $\lambda=m/\ell$，下垂部分质量为
$m_h=mx/\ell$，它随圆柱放绳以速率 $v=R\omega$
向下运动。缠绕部分虽质心位于轴上，其质量仍以角速度
$\omega$ 随圆柱转动，因此旋转惯量约为
$I_w=(m-m_h)R^2$；均匀圆柱惯量
$I_c=\frac12MR^2$。
将下垂绳的平动动能也用 $\omega$ 表示：
\[
K=\tfrac12 I_c\omega^2+
\tfrac12(m-m_h)R^2\omega^2+
\tfrac12m_hR^2\omega^2
=\tfrac12(\tfrac12M+m)R^2\omega^2.
\]
下垂部分对轴的重力力矩为
$\tau_g=m_hgR$；当绳从缠绕部分转换为下垂部分时，
上述总等效惯量不变。由角动量平衡：
\[
(\tfrac12M+m)R^2\beta
=\frac{mx}{\ell}gR.
\]
因此
\[
\boxed{\beta(x)=\frac{2mgx}{R\ell(M+2m)}}.
\]
\excheck 当 $x=0$ 时 $\beta=0$；若漏算绳的缠绕部分
转动惯量或下垂部分平动动能，都不能得到原书的分母 $M+2m$。

\exproblem{252}{39--40}
质量 $m$、半径 $R$ 的均匀实心球沿倾角为
$\alpha$ 的斜面无滑动滚下。求：
(a) 保证不打滑的最小静摩擦系数；
(b) 从静止起运动 $t$ 秒后的球的总动能。
\exsolution
令质心沿斜面向下加速度为 $a$，
接触静摩擦力大小为 $f$，方向沿斜面向上。
牛顿第二定律及绕球心的转动方程：
\[
mg\sin\alpha-f=ma,\quad
fR=I_C\beta,\quad
I_C=\tfrac25mR^2,\quad a=R\beta.
\]
由后两式 $f=\frac25ma$，
代入第一式得
$a=\frac57g\sin\alpha$，
$f=\frac27mg\sin\alpha$。
支持力 $N=mg\cos\alpha$，
要求 $f\le kN$，故
\[
\boxed{k\ge\frac27\tan\alpha}.
\]
从静止开始 $v=at$，
总动能为质心平动加自转：
\[
K=\frac12mv^2+\frac12I_C(v/R)^2
=\frac7{10}mv^2
=\boxed{\frac5{14}mg^2t^2\sin^2\alpha}.
\]
\excheck 静摩擦阻止滑动但接触点瞬时静止，
在理想静止斜面上它不作机械功；
总动能增加来自重力势能下降。

\exproblem{253}{40}
一根质量 $m=8.0\,\mathrm{kg}$、半径
$R=1.3\,\mathrm{cm}$ 的均匀实心圆柱
由两根对称细绳悬挂，绳端固定于天花板，
圆柱从静止开始因重力向下放绳（图1.60）。
求：(a) 每根绳的张力及圆柱角加速度；
(b) 重力瞬时功率 $P(t)$。
\exsolution
设质心向下加速度为 $a$、
每根绳张力为 $T$。两绳无滑动放开，
角加速度 $\beta=a/R$。
力平衡与转动方程为
\[
mg-2T=ma,\qquad
2TR=I_C\beta=\tfrac12mR^2(a/R).
\]
第二式给出 $2T=\frac12ma$，
代回得 $a=\frac23g$。
所以
\[
\boxed{T=\frac16mg\simeq13.1\,\mathrm N},
\qquad
\boxed{\beta=\frac{2g}{3R}\simeq5.03\times10^2\,
\mathrm{rad/s^2}}.
\]
从静止开始质心速度为 $v=at=(2g/3)t$，
重力功率等于重力与速度点积：
\[
\boxed{P_g(t)=mgv=\frac23mg^2t}.
\]
\excheck 对两根绳应使用\emph{总张力 $2T$}列质心方程，
不能误用单根张力。

\exproblem{254}{40}
一根均匀实心圆柱质量为 $m$，
两端细绳对称缠绕于圆柱且自由端固定在电梯天花板。
电梯从静止开始以恒定加速度 $a_0$ \emph{向上}运动。
求：(a) 圆柱相对电梯的加速度；
(b) 圆柱通过两绳对天花板的总作用力大小。
\exsolution
以电梯为参考系，并约定\emph{向下}为圆柱相对加速度
$a'$ 的正方向。因电梯向上加速，
电梯系的等效重力场为
$g_{\mathrm{eff}}=g+a_0$，而不是 $g-a_0$。
设每根绳张力为 $T$，圆柱实心惯量
$I_C=\frac12mR^2$。方程为
\[
m(g+a_0)-2T=ma',\qquad
2TR=I_C\frac{a'}R.
\]
于是 $2T=\frac12ma'$，
从而 $m(g+a_0)=\frac32ma'$：
\[
\boxed{a'=\frac23(g+a_0)\quad\text{（相对电梯向下）}},
\]
\[
\boxed{F_{\mathrm{ceiling}}=2T=\frac13m(g+a_0)
\quad\text{（对天花板向下）}}.
\]
\excheck 本版书末答案将两式写为 $g-a_0$，
与题面明确的“电梯向上加速”不相容；它对应
\emph{电梯向下加速}的另一模型。
本稿按原题方向列受力，并将差异列入勘误记录。

\exproblem{255}{40}
质量 $m=200\,\mathrm g$ 的线轴位于倾角
$\alpha=30^\circ$ 的光滑斜面上，
其细绳自由端固定于上方墙壁（原图1.61）。
线轴对自转轴的转动惯量为
$I=0.45\,\mathrm{g\,m^2}$，
绕绳层半径 $r=3.0\,\mathrm{cm}$。
绳无滑动地绕线轴放开，求线轴中心加速度。
\exsolution
沿斜面向下为正，设线轴中心加速度为 $a$，
绳张力为 $T$，角加速度为 $\beta$。
细绳与线轴接触处相对固定线端无滑动，
因此 $a=r\beta$。
沿斜面平动及绕质心转动方程分别为
\[
mg\sin\alpha-T=ma,\qquad
Tr=I\beta=I\frac ar.
\]
从而 $T=Ia/r^2$，得到
\[
\boxed{a=\frac{g\sin\alpha}{1+I/(mr^2)}}.
\]
按题面\emph{明确印出的单位}
$0.45\,\mathrm{g\,m^2}=4.5\times10^{-4}
\,\mathrm{kg\,m^2}$，代入
$I/(mr^2)=2.50$：
\[
\boxed{a\simeq\frac{9.81\cdot0.5}{3.50}
\simeq1.40\,\mathrm{m/s^2}}.
\]
\excheck 底本书末给约 $1.6\,\mathrm{m/s^2}$，
但题面数据和所印公式给出约 $1.4\,\mathrm{m/s^2}$。
这属于源书数值冲突，不能通过擅改惯量单位来掩盖。

\exproblem{256}{40}
质量为 $m$ 的均匀实心圆柱放在两块平行、水平的
粗糙窄支撑板上（图1.62）。
细绳缠绕在圆柱表面，自由端受恒定竖直向下拉力 $F$。
圆柱与支撑板之间的静摩擦系数为 $k$。
求使圆柱仍能在板上无滑动滚动的最大拉力
$F_{\max}$ 及此时圆柱轴线的最大平动加速度。
\exsolution
由图中的细绳下拉，圆柱作平面滚动；
沿滚动方向只有支撑板施加的水平静摩擦合力 $f$，
所以 $f=ma$。
关于圆柱质心的转动方程（下拉力和摩擦力矩相反）：
\[
(F-f)R=I_C\beta,\qquad
I_C=\tfrac12mR^2,\quad \beta=\frac aR.
\]
故 $F-f=\frac12ma$，
结合 $f=ma$ 可得
\[
\boxed{a=\frac{2F}{3m}},\qquad
f=\frac{2F}{3}.
\]
竖直方向支撑板共同支持
$N=mg+F$；在无滑动极限有
$f_{\max}=kN$。因此
\[
\frac{2F_{\max}}{3}=k(mg+F_{\max}),
\qquad
\boxed{F_{\max}=\frac{3kmg}{2-3k}},
\]
\[
\boxed{a_{\max}=\frac{2kg}{2-3k}}.
\]
\excheck 这两个有限上界只适用于 $0\le k<2/3$。
在理想库仑静摩擦模型且 $k\ge2/3$ 时，
上述不滑动条件对任意有限 $F$ 均可满足，
题目就没有有限的 $F_{\max}$。
原书第256题末尾误出现“rolling down the inclined plane”
字样，但图1.62明确是水平支撑板；
本稿采用图示装置，指出文字拼接歧义。

\exproblem{257}{40--41}
线轴质量 $m$，外半径 $R$、绕线层半径 $r$，
绕自转轴的转动惯量
$I=\gamma mR^2$。
线轴在粗糙\emph{水平}面上无滑动滚动。
按原图1.63，轻绳沿内层的上侧切向受恒力
$F$，拉力与水平正向 $x$ 夹角为 $\alpha$。
求：(a) 轴心加速度的 $x$ 分量；
(b) 从静止开始前 $t$ 秒拉力做功。
\exsolution
取轴心向右加速度为 $a$，
绕轴角加速度取逆时针为正，接触点无滑动意味着
$\beta=-a/R$。
记地面水平摩擦力为 $f$（以右为正），
平动方程为
$m a=F\cos\alpha+f$。
按原书图1.63的上侧出线方向，拉力
对轴心产生正向力矩 $Fr$。
转动方程和无滑动约束为
\[
I\beta=Fr+fR,\qquad \beta=-a/R.
\]
因此
\[
f=-Fr/R-Ia/R^2.
\]
与平动方程 $ma=F\cos\alpha+f$ 联立：
\[
(m+I/R^2)a=F(\cos\alpha-r/R).
\]
最终
\[
\boxed{a_x=\frac{F(\cos\alpha-r/R)}{m(1+\gamma)}}.
\]
由无滑动，地面静摩擦作用点相对地面瞬时静止，
故它不作功；刚体从静止开始总动能
\[
K=\tfrac12mv_C^2+\tfrac12I(v_C/R)^2
=\tfrac12m(1+\gamma)(a_xt)^2.
\]
因此拉力在前 $t$ 秒的功
\[
\boxed{W_F=\frac{F^2t^2(\cos\alpha-r/R)^2}
{2m(1+\gamma)}}.
\]
\excheck 当 $\cos\alpha=r/R$ 时
加速度和总功均为零（理想瞬态静止）。
两个方向的受力取向须与图1.63的绕绳方向对应。

\exproblem{258}{41}
原图1.64中，两根相同的均匀实心圆柱
每根质量为 $m$、半径为 $R$，
以两根对称、轻而不可伸长的细绳构成上下联动放绳装置。
上圆柱轴固定且无摩擦，下圆柱自由悬挂、下落时
两根绳在上下圆柱处均无滑动。
求运动时\emph{每根}细绳的张力 $T$。
\exsolution
上下两圆柱的自转惯量都为
$I=\frac12mR^2$，两根绳各张力为 $T$。
对上圆柱、下圆柱分别列转动方程：
\[
I\beta_1=2TR,\qquad I\beta_2=2TR,
\quad\Longrightarrow\quad
\beta_1=\beta_2=\frac{4T}{mR}.
\]
依图示放绳的两级滚动约束，
下圆柱质心向下加速度
$a=R(\beta_1+\beta_2)=8T/m$。
对下圆柱的质心列竖直方程：
\[
mg-2T=ma=8T.
\]
因而
\[
\boxed{T=\frac{mg}{10}},
\qquad
\boxed{a=\frac45g}.
\]
\excheck $T$ 是每根绳的张力，不是两绳总张力；
若误认为上圆柱固定不转，会改变放绳约束并得出错误结果。

\exproblem{259}{41}
原图1.65为复合卷筒装置：重物 $A$ 质量为 $m$，
可升降并转动的滑轮 $B$ 质量为 $M$，
关于其自转轴的惯量为 $I$。
两层卷绳半径分别为 $R$ 和 $2R$，
其中小半径绳一端固定于天花板，
大半径绳连接重物 $A$。从静止释放，
求重物 $A$ 向下的加速度。
\exsolution
设滑轮轴向下位移为 $y$，旋转角为 $\varphi$。
小半径绳上端固定且无滑动，
故 $\dot y=R\dot\varphi$。
大半径的绳从同向另一层放开，使
重物 $A$ 的下降速度满足
$\dot z=\dot y+2R\dot\varphi=3\dot y$。
因此 $\ddot z=3\ddot y$，
两层无滑动约束已将系统约化为一个自由度。
总动能为
\[
K=\tfrac12M\dot y^2+
\tfrac12I(\dot y/R)^2+
\tfrac12m(3\dot y)^2
=\tfrac12(M+I/R^2+9m)\dot y^2.
\]
当滑轮轴下降 $\mathrm dy$ 时，
重物下降 $3\mathrm dy$，重力作功
$\mathrm dW_g=(M+3m)g\,\mathrm dy$。
由拉格朗日方程
\[
(M+I/R^2+9m)\ddot y=(M+3m)g.
\]
故重物加速度
\[
\boxed{a_A=3\ddot y
=\frac{3g(M+3m)}{M+9m+I/R^2}}
\quad\text{（向下）}.
\]
\excheck 必须区分滑轮轴的加速度 $\ddot y$
与重物的加速度 $3\ddot y$；
此式与原书答案一致。

\exproblem{260}{41}
质量为 $m_1$ 的均匀实心圆柱 $A$
可绕固定于底座 $B$ 的水平轴无摩擦转动，
底座质量为 $m_2$ 且能在光滑水平面上平移（图1.66）。
轻绳绕在圆柱表面，水平自由端 $K$ 受恒力 $F$ 向右拉。
从静止开始，求：(a) $K$ 点加速度；
(b) 系统在前 $t$ 秒获得的总动能。
\exsolution
整个底座和圆柱的唯一水平外力为 $F$，
故底座和圆柱轴心共同的水平加速度为
\[
a_C=\frac{F}{m_1+m_2}.
\]
圆柱绕自身轴受恒定力矩 $FR$，
角加速度大小
$\beta=FR/I=2F/(m_1R)$，
因为 $I=\frac12m_1R^2$。
按图中放绳切向方向，
端点 $K$ 的平动加速度为底座平动
加圆柱边缘的切向相对加速度：
\[
a_K=a_C+R\beta
=\frac F{m_1+m_2}+\frac{2F}{m_1}.
\]
因此
\[
\boxed{a_K=\frac{F(3m_1+2m_2)}{m_1(m_1+m_2)}}.
\]
由于从静止开始、$a_K$ 恒定，
端点位移为 $s_K=\frac12a_Kt^2$。
地面光滑、轴无耗散，外力做功全变成总动能：
\[
\boxed{K_{\rm system}(t)=Fs_K
=\frac{F^2t^2(3m_1+2m_2)}
{2m_1(m_1+m_2)}}.
\]
\excheck $K$ 的功应按\emph{绳端位移}计算，
不能只用底座的位移，否则漏掉圆柱自转动能。
% === END 05-rigid-body-251-260.tex ===

% === BEGIN 05-rigid-body-261-270.tex blob 9c1af98e90d63839db97f5ba7de366e491c1d931 ===
% Irodov, user-provided PDF, original pp.41--43, PDF pp.50--52.
% Problem 264's incline slopes DOWN as drawn in Fig. 1.67, not uphill.

\exproblem{261}{41}
质量为 $m_1$ 的木板置于光滑水平面上，
质量为 $m_2$ 的均匀实心球置于板上。
对木板施加恒定水平力 $F$，
假设球与板之间始终无相对滑动。
求木板与球心的水平加速度 $a_1,a_2$。
\exsolution
设板向右加速，静摩擦力 $f$ 使球心向右加速并
使球绕自身中心逆时针自转。对球有
\[
m_2a_2=f,\qquad I_C\beta=fR,\qquad
I_C=\frac25m_2R^2.
\]
球与板接触处无相对切向加速度，
因此
$a_1=a_2+R\beta$。
消去 $\beta,f$：
\[
a_1=a_2+\frac{fR^2}{I_C}
=a_2+\frac52a_2=\frac72a_2.
\]
木板的水平运动方程
$F-f=m_1a_1$，代入 $f=m_2a_2=(2/7)m_2a_1$：
\[
\boxed{a_1=\frac{F}{m_1+\frac27m_2}},\qquad
\boxed{a_2=\frac27a_1}.
\]
\excheck 球心比木板加速得慢，
但无滑动条件约束的是球体接触点的速度，
不是两者质心速度相等。

\exproblem{262}{41--42}
均匀实心圆柱质量为 $m$、半径为 $R$，
最初仅绕自身轴以角速度 $\omega_0$ 自转、质心静止，
随后被放到水平粗糙面上释放。
摩擦系数为 $k$。求：
(a) 从释放到开始纯滚动的滑动持续时间；
(b) 滑动摩擦对圆柱所做\emph{总功}。
\exsolution
滑动期间地面摩擦力大小恒为 $f=kmg$。
其使圆柱获得质心速度，$a_C=kg$，
同时施加阻转力矩，
\[
|\beta|=\frac{fR}{I_C}=\frac{2kg}{R},
\qquad I_C=\frac12mR^2.
\]
以摩擦使质心前进的方向为正，
\[
v_C(t)=kgt,\qquad
\omega(t)=\omega_0-\frac{2kg}{R}t.
\]
首次纯滚动条件为 $v_C=R\omega$：
\[
kgt=\omega_0R-2kgt
\Longrightarrow
\boxed{t_{\mathrm{slide}}=\frac{\omega_0R}{3kg}}.
\]
此时 $v_f=\omega_0R/3,\ \omega_f=\omega_0/3$。
摩擦对整个圆柱的功等于总动能变化：
\[
W_f=\left[\frac12mv_f^2+\frac12I_C\omega_f^2\right]
-\frac12I_C\omega_0^2
=\left(\frac1{18}+\frac1{36}-\frac14\right)
mR^2\omega_0^2.
\]
故
\[
\boxed{W_f=-\frac16mR^2\omega_0^2}.
\]
\excheck 在滑动阶段存在机械能耗散；
接触点速率由不为零逐渐降至零，之后进入纯滚动。

\exproblem{263}{42}
一颗均匀实心小球半径为 $r$，
在半径为 $R$ 的固定大球表面从顶端几乎静止开始
无滑动滚下。
求小球恰离开大球表面瞬间的自转角速度大小。
\exsolution
小球质心绕大球心移动的圆半径为 $L=R+r$。
设质心的径矢与竖直向上方向夹角为 $\theta$，
质心速率为 $v$，小球自转角速度
$\omega=v/r$。
由能量守恒：
\[
mgL(1-\cos\theta)
=\frac12mv^2+\frac12\left(\frac25mr^2\right)\frac{v^2}{r^2}
=\frac7{10}mv^2.
\]
故
$v^2=\frac{10}{7}gL(1-\cos\theta)$。
径向向心方向列牛顿第二定律：
\[
mg\cos\theta-N=\frac{mv^2}{L}.
\]
离面瞬间 $N=0$，联立得到
$\cos\theta=10/17$，
$v^2=(10/17)gL$。
从而
\[
\boxed{\omega_{\mathrm{detach}}=
\frac vr=\sqrt{\frac{10g(R+r)}{17r^2}}}.
\]
\excheck 半径为 $R+r$ 的圆是\emph{球心轨迹}，
滚动约束的自转半径则为 $r$，不能把二者混用。

\exproblem{264}{42}
半径 $R=15\,\mathrm{cm}$ 的均匀实心圆柱
原来以速率 $v_0$ 在水平面上向右纯滚动，
其前方路面在拐角处\emph{向下}折成
倾角 $\alpha=30^\circ$ 的斜面（原图1.67）。
求使圆柱进入斜面段时\emph{不腾空跳离拐角}的
最大 $v_0$，并假设运动中无滑动。
\exsolution
图1.67的斜面向右\emph{下倾}。
圆柱过拐角的短阶段以拐角 $O$ 为瞬时接触点，
圆心沿半径 $R$ 的圆弧从竖直上方转过
$\theta\in[0,\alpha]$。
圆柱绕接触点的转动惯量是
$I_O=I_C+mR^2=\frac32mR^2$。
由于纯滚动 $v=R\dot\theta$，
其总动能 $K=\frac12I_O\dot\theta^2
=\frac34mv^2$。
圆心在下坡阶段降低
$R(1-\cos\theta)$，能量守恒：
\[
\frac34mv^2
=\frac34mv_0^2+mgR(1-\cos\theta),
\quad
v^2=v_0^2+\frac43gR(1-\cos\theta).
\]
沿指向拐角的径向列式：
$mg\cos\theta-N=mv^2/R$，
其中 $N$ 为拐角向外的支持力。
要一直保持接触须 $N\ge0$，即
\[
\frac{N}{mg}=\frac73\cos\theta-\frac43
-\frac{v_0^2}{gR}\ge0.
\]
由于 $\cos\theta$ 在区间中递减，
最严条件发生在 $\theta=\alpha$，于是
\[
\boxed{v_{0,\max}=
\sqrt{\frac{gR}{3}(7\cos\alpha-4)}
\simeq1.00\,\mathrm{m/s}}.
\]
\excheck 根号内须为非负，即
$\cos\alpha\ge4/7$；本题 $\alpha=30^\circ$ 满足。
如果误把原图的向下斜坡当成向上斜坡，
得到的离面条件将完全不同。

\exproblem{265}{42}
一根均匀薄刚性圆环质量为 $m$、半径为 $R$，
其内缘固定着质量同样为 $m$ 的小物体 $A$。
圆环在水平面上无滑动滚动。
每当 $A$ 到达环内最低位置时，
圆环中心速率为 $v_0$（图1.68）。
求圆环始终不跳离地面的 $v_0$ 条件。
\exsolution
令 $\theta=0$ 表示小物体位于环心正下方，
圆环滚动时其中心速率为 $v=R\dot\theta$。
环本身的平动与转动动能之和为 $mv^2$；
小物体绕地面接触点作摆线运动，其速率平方为
$v_A^2=2v^2(1-\cos\theta)$，
故全系统动能
\[
K=mv^2+\frac12m v_A^2
=mv^2(2-\cos\theta).
\]
小物体重力势能相对最低点的增量为
$mgR(1-\cos\theta)$。能量守恒：
\[
v^2(2-\cos\theta)+gR(1-\cos\theta)=v_0^2.
\]
环和小物体总质心高度为
$y_C=R-\frac12R\cos\theta$，
质量合计 $2m$。
地面对系统的支持力
$N=2mg+2m\,\ddot y_C$；
代入上式与其关于时间的导数可确定
$N(\theta)$。在足以完成整圈的运动中，
地面支持力的最小值出现在
$A$ 位于最高位置 $\theta=\pi$。
此时 $\dot\theta^2=v^2/R^2$、
$\ddot y_C=-R\dot\theta^2/2$，于是
$N_{\mathrm{top}}=2mg-mv_{\mathrm{top}}^2/R$。
能量式在 $\theta=\pi$ 给出
$v_{\mathrm{top}}^2=(v_0^2-2gR)/3$。
要求 $N_{\mathrm{top}}\ge0$，
即
\[
\frac{v_0^2-2gR}{3}\le2gR
\quad\Longrightarrow\quad
\boxed{v_0\le\sqrt{8gR}}.
\]
\excheck 若 $v_0$ 太小，小物体不会到达圆环顶端，
但系统仍可能往复滚动而不跳起；
式中给的是无跳离地面的上限，原书亦如此给出。

\exproblem{266}{42}
一辆拖拉机以相对于地面恒速 $v$ 前进，
整条履带总质量为 $m$，如图1.69。
在无滑动滚动条件下，求整条履带相对地面的总动能。
\exsolution
在理想履带模型中，上方直线段沿前进方向
相对地面速率为 $2v$；底部接触地面的直线段
瞬时速率为零。忽略弯曲端部的质量，
两段履带质量分别为 $m/2$。
因此
\[
K=\frac12\left(\frac m2\right)(2v)^2
+\frac12\left(\frac m2\right)0^2
=\boxed{mv^2}.
\]
\excheck 若端部弯曲段质量不可忽略，应将沿履带各质元
的实际速度积分，此处采用原题图的理想化直线段近似。

\exproblem{267}{42}
质量为 $m$、半径为 $r$ 的均匀实心球
绕水平杆 $OA$ 转动并在水平面上无滑动滚动
（原图1.70）。球心沿距竖直旋转中心 $O$
为 $R$ 的圆轨道以速率 $v$ 前进。
求球的总动能。
\exsolution
球心平动动能为 $mv^2/2$。
运动中球同时绕水平径向轴滚动，角速度大小
$\omega_r=v/r$；球心的轨道绕竖直方向以
$\Omega=v/R$ 转动。两转动角速度矢量互相垂直。
实心球各直径轴惯量相同，均为
$I_C=\frac25mr^2$，故球自转动能为
\[
K_{\rm rot}=\frac12I_C(\omega_r^2+\Omega^2)
=\frac15m v^2\left(1+\frac{r^2}{R^2}\right).
\]
总动能
\[
\boxed{K=\frac7{10}mv^2\left(1+\frac27\frac{r^2}{R^2}\right)}.
\]
\excheck 若遗漏球心绕中心 $O$ 的方位转动角速度
$\Omega$，将丢失与 $r^2/R^2$ 成比例的附加动能。

\exproblem{268}{42--43}
证明：一个质量为 $m$、质心相对于固定旋转轴
的垂直位置矢量为 $\bm R_C$ 的刚体，
在以恒角速度 $\bm\omega$ 旋转的参考系中，
全部质元所受的离心惯性力合力和科里奥利惯性力合力分别为
\[
\bm F_{\rm cf}=m\omega^2\bm R_C,\qquad
\bm F_{\rm Cor}=2m(\bm v_C'\times\bm\omega),
\]
其中 $\bm v_C'$ 为质心在旋转系中的速度。
\exsolution
把刚体分成质元 $\mathrm dm$。
各质元离心惯性力
$\mathrm d\bm F_{\rm cf}
=-\bm\omega\times(\bm\omega\times\bm r_\perp)
\,\mathrm dm$。
由于 $\bm r_\perp$ 垂直于 $\bm\omega$，
双叉乘恒等式给出
$-\bm\omega\times(\bm\omega\times\bm r_\perp)
=\omega^2\bm r_\perp$。
对全体质量元积分：
\[
\bm F_{\rm cf}
=\omega^2\int\bm r_\perp\,\mathrm dm
=m\omega^2\bm R_C.
\]
各质元科里奥利惯性力为
$\mathrm d\bm F_{\rm Cor}
=-2\bm\omega\times\bm v_i'\,\mathrm dm$。
利用积分与叉积的线性性及质心速度定义：
\[
\bm F_{\rm Cor}=-2\bm\omega\times
\int\bm v_i'\,\mathrm dm
=-2m\bm\omega\times\bm v_C'
=\boxed{2m(\bm v_C'\times\bm\omega)}.
\]
\excheck 这些是\emph{合力}恒等式，不表示各质元受惯性力
均相同；求惯性力矩仍需考虑物体内部质量分布。

\exproblem{269}{43}
长为 $\ell$、质量为 $m$ 的均匀细杆中点 $C$
固定于旋转轴 $OO'$ 上，杆与轴夹角为 $\theta$
（原图1.71）。轴以恒定角速度 $\omega$ 转动。
求在随轴旋转的参考系中，全部离心惯性力
对杆中点 $C$ 的\emph{合力矩大小}。
\exsolution
令杆上质元相对 $C$ 的位置为
$\bm r=s\bm e$，其中
$s\in[-\ell/2,\ell/2]$、$\bm e$ 是杆轴单位矢量，
与竖直旋转轴夹角为 $\theta$。
离心惯性力为
$\mathrm d\bm F=\omega^2(s\sin\theta)
\bm e_\perp\,\mathrm dm$，
其中 $\bm e_\perp$ 是杆方向的水平投影的单位矢量。
每个微元关于 $C$ 的力矩模为
\[
\mathrm dN
=\omega^2s^2\sin\theta\cos\theta\,\mathrm dm,
\]
所有微元力矩同向。因均匀细杆
$\int s^2\,\mathrm dm=m\ell^2/12$，
积分得
\[
\boxed{N=\frac{m\omega^2\ell^2}{12}
\sin\theta\cos\theta
=\frac{m\omega^2\ell^2}{24}\sin2\theta}.
\]
\excheck 这也可由刚体惯量张量的陀螺力矩公式得到。
本版书末第269题印作 $\frac12m\omega^2\ell^2\sin2\theta$，
与上述积分及同一原书第282(c)题的
$\frac1{24}m\omega^2\ell^2\sin2\theta$ 不一致；
本稿保留可复核的连续体积分结论并标为勘误。

\exproblem{270}{43}
一圆锥摆由长为 $\ell$、质量为 $m$ 的均匀细杆构成，
上端通过铰链与固定点连接。
细杆绕竖直轴以恒角速度 $\omega$ 匀速旋转，
求其与竖直方向的平衡夹角 $\theta$。
\exsolution
在随杆匀速旋转的参考系中，
杆处于相对静止状态。
选铰点为力矩中心，铰接反力无力矩。
重力作用在杆中点，其使杆向竖直方向回落的力矩大小
\[
N_g=mg\frac{\ell}{2}\sin\theta.
\]
取距铰点沿杆距离为 $s$ 的质量微元，
其到竖直轴距离为 $s\sin\theta$，
离心力大小 $\omega^2s\sin\theta\,\mathrm dm$，
对铰点的力矩大小是
$\omega^2s^2\sin\theta\cos\theta\,\mathrm dm$。
积分且 $\int s^2\,\mathrm dm=m\ell^2/3$：
\[
N_{\rm cf}=\frac13m\omega^2\ell^2
\sin\theta\cos\theta.
\]
令力矩平衡 $N_g=N_{\rm cf}$，
对非竖直解约去 $\sin\theta$：
\[
\boxed{\cos\theta=\frac{3g}{2\omega^2\ell}}.
\]
\excheck 要存在非竖直圆锥摆，
必须 $\omega^2\ell\ge3g/2$；
低于该阈值时仅有 $\theta=0$ 的竖直平衡。
% === END 05-rigid-body-261-270.tex ===

% === BEGIN 05-rigid-body-271-280.tex blob 5e23ba432786abce2d37b8cc47b3a2cfc377244e ===
% Irodov, user-supplied PDF, problems 271--280, printed pp.43--45
% (PDF pp.52--54). Source diagrams 1.71--1.72 and answer sheet checked.

\exproblem{271}{43}
边长为 $a$ 的均匀立方体放在摩擦系数为 $k$ 的
水平面上，获得一初始水平速度后减速至停止。
解释它在滑行减速时关于\emph{固定于水平地面且垂直于运动方向的轴}
的角动量为何逐渐减少并消失，
并求地面支持力合力作用线相对重力作用线的水平偏移
$\Delta x$。
\exsolution
设立方体质量为 $m$，运动方向取为 $+x$，
立方体质心离地高度为 $a/2$。
滑动摩擦力大小为 $f=kmg$，指向 $-x$。
摩擦作用于底面，关于\emph{经过质心}的横向水平轴产生大小
\[
N_f=f(a/2)=kmg(a/2)
\]
的倾覆力矩。若立方体只平动而不倾转，
总转矩须为零。重力作用线穿过质心，
而地面总支持力 $N=mg$ 可以在底面上偏离几何中心，
对该轴产生反向力矩 $N\,\Delta x$。
因此
\[
mg\,\Delta x=kmg\,\frac a2
\quad\Longrightarrow\quad
\boxed{\Delta x=\frac12ka}.
\]
由于地面摩擦使质心减速，立方体关于\emph{固定于地面}
的横向水平轴所具有的平动轨道角动量逐渐减少至零。
支持力作用线偏移只保证关于\emph{质心}的总力矩为零、
不发生自转，并不能阻止平动角动量的变化。
\excheck 压力中心的偏移源自底面压力分布的不对称；
若 $\Delta x>a/2$，即 $k>1$，上述不倾转状态
不再自洽，应考虑倾翻。

\exproblem{272}{43}
一根长为 $\ell$、质量为 $M$ 的均匀水平细杆 $AB$
可绕通过端点 $A$ 的竖直固定轴无摩擦旋转，
最初角速度为 $\omega_0$。
质量为 $m$ 的小套环从轴心 $A$ 开始无摩擦地
沿杆滑向另一端 $B$。
求套环到达 $B$ 时相对杆的径向滑动速率 $v'$。
\exsolution
杆对端点的转动惯量为
$I_0=M\ell^2/3$。
套环离轴距离为 $r$ 时，总角动量
$L_z=(I_0+mr^2)\omega$。
外力对固定竖轴的力矩为零，故
\[
(I_0+mr^2)\omega=I_0\omega_0.
\]
套环在 $r=\ell$ 时
$\omega_f=I_0\omega_0/(I_0+m\ell^2)$。
无摩擦且轴不作功，机械能守恒：
\[
\tfrac12I_0\omega_0^2
=\tfrac12(I_0+m\ell^2)\omega_f^2
+\tfrac12m{v'}^2.
\]
代入 $\omega_f$：
\[
m{v'}^2
=I_0\omega_0^2
\left[1-\frac{I_0}{I_0+m\ell^2}\right]
=\frac{I_0m\ell^2\omega_0^2}{I_0+m\ell^2}.
\]
因此
\[
\boxed{v'=\frac{\omega_0\ell}{\sqrt{1+3m/M}}}.
\]
\excheck $v'$ 是套环沿杆的\emph{相对径向速率}，
不是套环对地总速率，后者还含方位速度分量。

\exproblem{273}{43}
质量为 $m=5.0\,\mathrm{kg}$、
长度为 $\ell=90\,\mathrm{cm}$ 的均匀细杆
在光滑水平面上静止，一端突然受到
大小为 $J=3.0\,\mathrm{N\,s}$、
垂直于杆的水平冲量。
求随后的运动中，杆一半对另一半的内力大小。
\exsolution
冲量使杆的质心获得平动速度
$v_C=J/m$；关于质心的冲量矩为
$J\ell/2$。
均匀细杆对质心的转动惯量为
$I_C=m\ell^2/12$，所以碰后角速度
\[
\omega=\frac{J\ell/2}{I_C}
=\frac{6J}{m\ell}.
\]
冲量结束后外部水平合力为零，
质心匀速运动、角速度恒定。
每半杆的质心距整杆质心 $\ell/4$，
做圆周运动所需向心加速度为
$\omega^2\ell/4$。
半杆质量 $m/2$，沿杆指向整体质心的
内力大小因此为
\[
F_{\rm internal}=\frac m2\cdot\frac{\omega^2\ell}{4}
=\frac{m\ell\omega^2}{8}.
\]
代入得
\[
\boxed{F=\frac{9J^2}{2m\ell}
=\frac{9(3.0)^2}{2(5.0)(0.90)}
=9.0\,\mathrm N}.
\]
\excheck 该内力是在\emph{碰撞结束后的转动阶段}
维持两半杆的向心加速度，
不是瞬时冲量 $J$ 除以未知碰撞时间得到的力。

\exproblem{274}{43}
质量为 $M$、边长为 $\ell$ 的均匀正方形薄板，
绕其一条竖直侧边形成的固定轴无摩擦转动。
质量为 $m$ 的小球以速度 $v$ 垂直薄板撞击其中心，
发生完全弹性碰撞。
求：(a) 小球碰后有符号速度 $v'$；
(b) 碰后薄板轴承受到的水平合力大小。
\exsolution
薄板对过边的竖直轴惯量
$I=\frac13M\ell^2$，
撞击点距轴 $b=\ell/2$。
设碰撞时小球对薄板冲量为 $J$，
则小球碰后速度 $v'=v-J/m$，
薄板角速度
\[
\omega=\frac{Jb}{I}.
\]
完全弹性的一维局部接触恢复条件为
$b\omega-v'=v$，即
$J(b^2/I+1/m)=2v$。
因 $b^2/I=3/(4M)$，
\[
J=\frac{2v}{1/m+3/(4M)}
=\frac{8mMv}{4M+3m}.
\]
于是
\[
\boxed{v'=v\frac{3m-4M}{3m+4M}}.
\]
碰后板质心距转轴 $\ell/2$，
它以角速度 $\omega$ 作圆周运动，
所需向心加速度为 $\omega^2\ell/2$，
因此轴的水平合反力大小为
\[
F=M\omega^2\frac\ell2,\qquad
\omega=\frac{J\ell/2}{M\ell^2/3}
=\frac{12mv}{\ell(4M+3m)}.
\]
故
\[
\boxed{F=\frac{72Mm^2v^2}{\ell(4M+3m)^2}
=\frac{8Mv^2}{\ell[1+4M/(3m)]^2}}.
\]
\excheck $v'$ 可为负，意味着小球反弹；
碰撞时角动量关于轴守恒，但整体线动量会受轴冲量，
不能直接令两物体的线动量守恒。

\exproblem{275}{43--44}
一根质量为 $M$、长度为 $\ell$ 的均匀细杆
上端悬挂在固定水平铰点，初始竖直静止。
质量为 $m$ 的子弹水平飞来击中杆下端并嵌入，
随后杆摆过最大角度 $\alpha$。设 $m\ll M$。
求：(a) 子弹飞来速率；
(b) 碰撞过程中“子弹+杆”系统\emph{线动量增量}
及其原因；(c) 为使碰撞过程中系统线动量不变，
子弹应击中距上端多远的点 $x$？
\exsolution
碰撞很短，关于铰点的冲量矩为零，
故子弹撞前角动量 $mv\ell$ 等于碰后杆子系统的
角动量 $I\omega$，其中
$I\simeq M\ell^2/3$，从而
$\omega\simeq3mv/(M\ell)$。
碰后摆起阶段机械能守恒。
最大角处杆的质心升高
$\frac\ell2(1-\cos\alpha)$，
忽略 $m/M$ 高阶修正，
\[
\frac12I\omega^2
=\frac12Mg\ell(1-\cos\alpha).
\]
联立求得
\[
\boxed{v\simeq\frac M m
\sqrt{\frac{2g\ell}{3}}\sin\frac\alpha2}.
\]
碰后线动量主要来自杆的质心运动：
$P_f\simeq M(\omega\ell/2)
=\frac32mv$，碰前为 $P_i=mv$。
故碰撞过程中的水平线动量增量
\[
\boxed{\Delta p\simeq\frac12mv
=\frac M2\sqrt{\frac{2g\ell}{3}}\sin\frac\alpha2
=M\sqrt{\frac{g\ell}{6}}\sin\frac\alpha2}.
\]
它由固定铰点施加的\emph{外部水平冲量}造成。
若击中距上端为 $x$ 的位置，
碰后 $\omega\simeq mvx/I$，
线动量 $P_f\simeq M(\ell/2)\omega$。
令 $P_f=mv$ 得
\[
\boxed{x=\frac{2I}{M\ell}\simeq\frac{2\ell}{3}}.
\]
\excheck 此处始终采用 $m\ll M$；
撞击阶段不用机械能守恒，摆起阶段才使用。

\exproblem{276}{44}
一根质量为 $M$、半径为 $R$ 的均匀水平圆盘
可绕过中心的固定竖轴自由旋转。
盘上径向导槽中有质量为 $m$ 的小物体，
以穿过轴心的轻绳相连。
初始物体位于盘边缘，整体以角速度 $\omega_0$ 转动。
随后缓慢拉绳使物体移到轴心。
求：(a) 最后角速度；(b) 拉绳外力做功。
\exsolution
盘的转动惯量
$I_d=\frac12MR^2$。
拉绳通过竖轴，不提供关于轴的外力矩，
故总角动量守恒：
\[
(I_d+mR^2)\omega_0=I_d\omega_f.
\]
由此
\[
\boxed{\omega_f=\left(1+\frac{2m}{M}\right)\omega_0}.
\]
缓慢操作使最终径向振动动能可忽略，
拉绳的外力功等于系统转动动能增量：
\[
W=\frac12I_d\omega_f^2
-\frac12(I_d+mR^2)\omega_0^2.
\]
将 $\omega_f$ 代入：
\[
\boxed{W=\frac12mR^2\omega_0^2
\left(1+\frac{2m}{M}\right)}.
\]
\excheck 角动量守恒但动能\emph{增加}，
多出的能量来自拉绳外力做功。

\exproblem{277}{44}
均匀水平圆盘质量为 $m_2$、半径为 $R$，
能绕中心固定竖轴无摩擦转动。
质量为 $m_1$ 的人站在盘边缘。
原先整体静止，此后人沿盘边行走，
相对于圆盘累计转过角度 $\varphi'$，
相对于盘面的切向速率为 $v'(t)$，最后停下。
忽略人的尺寸，求：
(a) 圆盘相对地面转过角度 $\varphi$；
(b) 人对圆盘施加的竖轴力矩 $N(t)$。
\exsolution
圆盘惯量
$I_d=\frac12m_2R^2$。
人在地面系中的角速度为
$\dot\varphi+\dot\varphi'$。
因整个系统关于竖轴无外力矩，初始角动量零：
\[
I_d\dot\varphi+
m_1R^2(\dot\varphi+\dot\varphi')=0.
\]
所以
$\dot\varphi=
-\frac{m_1R^2}{I_d+m_1R^2}\dot\varphi'$。
由零初角积分：
\[
\boxed{\varphi=-\frac{2m_1}{2m_1+m_2}\varphi'}.
\]
人对圆盘的力矩等于盘角动量变化率：
\[
N=I_d\ddot\varphi
=-\frac{I_d m_1R^2}{I_d+m_1R^2}\ddot\varphi'.
\]
因 $v'=R\dot\varphi'$，
\[
\boxed{N(t)=-\frac{m_1m_2R}{2m_1+m_2}
\frac{\mathrm dv'}{\mathrm dt}}.
\]
\excheck 人加速向前走时盘反向转动；
当人减速并最终相对盘停住时两者地面系角速度
也同时回到零，但盘可能已有非零角位移。

\exproblem{278}{44}
两个水平圆盘绕同一固定竖轴自由旋转，
惯量分别为 $I_1,I_2$，初始角速度为
$\omega_1,\omega_2$。
上盘落到下盘，最终因摩擦一起以同一角速度旋转。
求：(a) 最终角速度；(b) 此过程摩擦力做的总功。
\exsolution
轴无外力矩，故总角动量守恒：
\[
I_1\omega_1+I_2\omega_2=(I_1+I_2)\omega.
\]
因此
\[
\boxed{\omega=\frac{I_1\omega_1+I_2\omega_2}{I_1+I_2}}.
\]
两圆盘之间的滑动摩擦使动能损失，
其对两盘的\emph{合计做功}等于末初动能差：
\[
W_f=\frac12(I_1+I_2)\omega^2
-\frac12I_1\omega_1^2-\frac12I_2\omega_2^2.
\]
代入公共角速度并整理：
\[
\boxed{W_f=-\frac{I_1I_2}{2(I_1+I_2)}
(\omega_1-\omega_2)^2}.
\]
\excheck $W_f\le0$，且仅当
$\omega_1=\omega_2$ 时等于零；
角动量守恒不表示转动动能守恒。

\exproblem{279}{44}
一小圆片质量为 $m$，在光滑水平面上以速度 $v$
沿垂直于一根静止均匀细杆的方向运动，
随后与该杆的一端发生完全弹性碰撞。
细杆长为 $\ell$、质量为 $\eta m$，
其中 $\eta>0$。求碰后圆片速度 $v'$、
细杆角速度 $\omega$；
并求圆片碰后停下或反弹的质量比条件。
\exsolution
设沿入射方向的碰撞冲量大小为 $J$。
圆片碰后有符号速度
$v'=v-J/m$。
杆的质心平动速率
$V=J/(\eta m)$。
杆关于质心的惯量
$I=\eta m\ell^2/12$，
碰撞点距质心为 $\ell/2$，
故其角速度大小
\[
\omega=\frac{J(\ell/2)}{I}
=\frac{6J}{\eta m\ell}.
\]
杆端沿碰撞方向的速率为
$V+\omega\ell/2=4J/(\eta m)$。
弹性碰撞恢复系数为1，因而两接触点
碰后的相对分离速率等于碰前接近速率：
\[
\frac{4J}{\eta m}-\left(v-\frac Jm\right)=v.
\]
所以 $J=\frac{2\eta m v}{4+\eta}$，
\[
\boxed{v'=v\frac{4-\eta}{4+\eta}},
\qquad
\boxed{\omega=\frac{12v}{\ell(4+\eta)}}.
\]
\excheck 当 $\eta=4$，$v'=0$；
当 $\eta>4$，$v'<0$，圆片沿反方向弹回；
$0<\eta<4$ 时圆片仍向前运动。

\exproblem{280}{45}
图1.72的水平平台 $P$ 可绕固定竖轴 $OO'$
自由旋转；平台、电动机和配重关于此轴的惯量为 $I$。
电动机装有轻框架，框架支承一颗可沿自身轴 $BB'$
自由自转的均匀球 $A$；球关于任一直径的惯量为 $I_0$。
初始 $BB'$ 与 $OO'$ 重合，
球以自转角速度 $\omega_0$ 旋转，平台原先静止。
电动机慢慢将自转轴 $BB'$ 从竖直方向转过
$90^\circ$ 或 $180^\circ$，求：
(a) 两种转角所需电动机功；
(b) 轴转至 $90^\circ$ 后，为保持平台轴竖直
所需外力矩的大小。
\exsolution
设球自转轴与竖直方向的瞬时夹角为 $\theta$，
平台绕竖直轴的角速度为 $\Omega$。
由于球绕自身轴\emph{自由}自转，
平台自转对球附加的角速度
仅为\emph{垂直于}球自转轴的分量，
其大小为 $\Omega\sin\theta$。
因此关于竖直轴的总角动量为
\[
L_z=[I+I_0\sin^2\theta]\Omega
+I_0\omega_0\cos\theta.
\]
系统竖直轴上无外力矩，初始
$L_z=I_0\omega_0$，所以
\[
\Omega(\theta)=
\frac{I_0\omega_0(1-\cos\theta)}
{I+I_0\sin^2\theta}.
\]
转轴缓慢改变时忽略其瞬态动能，
相对于初始纯自转的附加动能为
\[
\Delta K=\frac12[I+I_0\sin^2\theta]\Omega^2
=\frac{I_0^2\omega_0^2(1-\cos\theta)^2}
{2[I+I_0\sin^2\theta]}.
\]
无耗散时该增量等于电动机做功，故
\[
\boxed{W_{90^\circ}=
\frac{I_0^2\omega_0^2}{2(I+I_0)}},\qquad
\boxed{W_{180^\circ}=\frac{2I_0^2\omega_0^2}{I}}.
\]
在 $\theta=90^\circ$ 时平台角速度为
$\Omega=I_0\omega_0/(I+I_0)$。
此时球的主自转角动量
$I_0\omega_0$ 位于水平面，
该矢量随平台以 $\Omega$ 绕竖直轴转动，
保持轴线方向所需外力矩大小等于其方向变化率：
\[
\boxed{N=I_0\omega_0\Omega
=\frac{I_0^2\omega_0^2}{I+I_0}}.
\]
\excheck 若误认为球与平台绕竖直轴\emph{刚性锁定}，
就会在 $\theta=180^\circ$ 得到错误的转动惯量；
原题明确说球可绕 $BB'$ 自由转动，
必须区分自转轴向和平面正交分量。
% === END 05-rigid-body-271-280.tex ===

% === BEGIN 05-rigid-body-281-289.tex blob 83c077ccaa5ad93c5408925a0f20c62a61fb6980 ===
% Irodov, user-supplied PDF, original printed pp.45--46, PDF pp.54--55.
% Rigid body: rotating bearings, gyroscopic precession, and rail loading.

\exproblem{281}{45}
一根质量 $m=1.40\,\mathrm{kg}$、长度
$\ell_0=1.00\,\mathrm m$ 的均匀细杆水平放置，
一端 $A$ 固定在竖直转轴 $OO'$ 的中点，
可随竖轴以角速度 $\omega$ 匀速转动。
轴上下两个轴承距离为
$\ell=55\,\mathrm{cm}$。
求当下轴承所受水平力为零时所需角速度，
并求此时上轴承所受水平力大小。
\exsolution
杆的质心距轴 $\ell_0/2$，需要水平向心力
$F_c=m\omega^2\ell_0/2$。
杆重力关于轴中点产生大小
$N_g=mg\ell_0/2$ 的倾覆力矩；
上下轴承的水平反力形成力偶以平衡它。
当下轴承的水平反力为零时，
上轴承的水平力 $F_u$ 必须同时满足
\[
F_u=F_c=\frac12m\omega^2\ell_0,
\qquad
F_u\frac{\ell}{2}=mg\frac{\ell_0}{2}.
\]
联立得到
\[
\boxed{\omega=\sqrt{\frac{2g}{\ell}}
\simeq5.97\,\mathrm{rad/s}\approx6.0\,\mathrm{rad/s}},
\]
\[
\boxed{F_u=\frac{mg\ell_0}{\ell}
\simeq25.0\,\mathrm N}.
\]
\excheck 杆固定在竖轴的\emph{中点}，
上轴承到固定点的力臂是 $\ell/2$，
不是全部轴承间距 $\ell$；
采用正确力矩臂后得到 $25\,\mathrm N$，
与底本书末答案一致。

\exproblem{282}{45}
长为 $\ell$、质量为 $m$ 的均匀细杆中点 $C$
固定于竖直轴 $OO'$，
杆与轴夹角为 $\theta$（原图1.71）。
系统绕竖轴以恒角速度 $\omega$ 旋转。
求：(a) 杆关于中点的角动量矢量大小、方向及其竖轴分量；
(b) 半周内角动量矢量增量的模；
(c) 轴承必须提供的外力矩大小。
\exsolution
杆对质心的转动惯量张量沿杆轴分量近似为零，
垂直杆轴的分量
$I_\perp=m\ell^2/12$。
记杆方向单位矢量为 $\bm e$，
竖直单位矢量为 $\bm k$，
转动角速度矢量 $\bm\omega=\omega\bm k$。
关于中点的角动量
\[
\bm L_C=I_\perp[
\bm\omega-(\bm\omega\cdot\bm e)\bm e]
=I_\perp\omega(\bm k-\cos\theta\,\bm e).
\]
其方向在杆与竖轴张成的平面内、垂直于杆，
大小和竖轴投影为
\[
\boxed{|\bm L_C|=\frac{m\ell^2\omega}{12}\sin\theta},
\qquad
\boxed{L_z=\frac{m\ell^2\omega}{12}\sin^2\theta}.
\]
随系统转过半周，角动量的水平分量反向，
竖直分量不变。水平分量大小为
$I_\perp\omega\sin\theta\cos\theta$，
故
\[
\boxed{|\Delta\bm L_C|
=2I_\perp\omega|\sin\theta\cos\theta|
=\frac{m\ell^2\omega}{12}|\sin2\theta|}.
\]
该水平分量以角速度 $\omega$ 旋转，
所以轴承需要的力矩
\[
\boxed{N=\omega|L_{\perp,z}|
=\frac{m\ell^2\omega^2}{24}|\sin2\theta|}.
\]
\excheck 这里 $L_{\perp,z}$ 表示角动量相对于
\emph{竖轴}的水平分量。
(c) 与本书第269题同一几何结构的连续积分应一致；
本题书末给 $\frac1{24}$，而第269题简答给 $\frac12$，
进一步表明后者存在系数印误。

\exproblem{283}{45}
一个陀螺质量为 $m=0.50\,\mathrm{kg}$，
自转轴与竖直夹角 $\theta=30^\circ$，
绕自转轴惯量 $I=2.0\,\mathrm{g\,m^2}$，
自转角速度 $\omega=350\,\mathrm{rad/s}$。
支点到质心的距离 $\ell=10\,\mathrm{cm}$。
在慢进动近似下求：(a) 进动角速度；
(b) 支点对陀螺的水平反力大小和方向。
\exsolution
自转主角动量模为 $L_s=I\omega$。
重力关于支点的力矩大小为
$N_g=mg\ell\sin\theta$。
慢进动时 $\bm L_s$ 以角速度 $\Omega$
绕竖直方向转动，
其矢量变化率大小为 $\Omega L_s\sin\theta$。
因此
\[
\Omega I\omega\sin\theta=mg\ell\sin\theta,
\qquad
\boxed{\Omega=\frac{mg\ell}{I\omega}
\simeq0.701\,\mathrm{rad/s}}.
\]
质心绕竖轴以半径 $\ell\sin\theta$
作匀速圆周运动，水平向心加速度为
$\Omega^2\ell\sin\theta$。
重力无水平分量，故支点水平反力为
\[
\boxed{F_h=m\Omega^2\ell\sin\theta
\simeq1.23\times10^{-2}\,\mathrm N},
\]
方向\emph{水平指向进动圆的旋转轴}。
\excheck 所求是慢进动模型下的支点反力，
不等于陀螺所受重力大小；
原书将其量级约记为 $10\,\mathrm{mN}$。

\exproblem{284}{46}
一陀螺由半径 $R=5.0\,\mathrm{cm}$ 的
均匀圆盘及长 $\ell=10\,\mathrm{cm}$ 的轻杆组成，
杆一端铰接在竖直向上以
$a_0=2.0\,\mathrm{m/s^2}$ 加速的电梯地板上（图1.73）。
圆盘以稳定角速度
$n=0.5$ 转/秒进动。
忽略摩擦及杆质量，求圆盘自身的角速度 $\omega$。
\exsolution
电梯参考系的等效重力加速度为
$g_{\rm eff}=g+a_0$。
均匀圆盘自转惯量
$I=\frac12mR^2$，主角动量
$L_s=I\omega$。
以 $\ell$ 为力臂，重力与惯性力合力矩
$mg_{\rm eff}\ell$ 引起慢进动，
其角速度 $\Omega=2\pi n$，故
\[
\Omega I\omega=mg_{\rm eff}\ell.
\]
消去圆盘质量得
\[
\boxed{\omega=
\frac{2(g+a_0)\ell}{R^2(2\pi n)}
=\frac{(g+a_0)\ell}{\pi nR^2}
\simeq3.01\times10^2\,\mathrm{rad/s}}.
\]
\excheck 电梯\emph{向上}加速使等效重力变大，
因此公式必须含 $g+a_0$；原书简答与此一致。

\exproblem{285}{46}
陀螺质量 $m=1.0\,\mathrm{kg}$，
自转轴转动惯量 $I=4.0\,\mathrm{g\,m^2}$、
自转角速度 $\omega=310\,\mathrm{rad/s}$。
支点固定在一水平加速的物块上，
物块加速度大小为 $a_0=1.0\,\mathrm{m/s^2}$。
支点到质心距离 $\ell=10\,\mathrm{cm}$。
求陀螺稳定慢进动角速度的模及进动轴方向。
\exsolution
在随支点加速的非惯性系中，
除重力 $m\bm g$ 外还有与支点水平加速方向相反的
惯性力 $-m\bm a_0$。
二者合成有效重力场
$\bm g_{\rm eff}=\bm g-\bm a_0$，
其大小
$g_{\rm eff}=\sqrt{g^2+a_0^2}$。
慢进动时
$\Omega I\omega=m\ell g_{\rm eff}$，
所以
\[
\boxed{\Omega=\frac{m\ell\sqrt{g^2+a_0^2}}{I\omega}
\simeq0.796\,\mathrm{rad/s}}.
\]
进动轴与有效重力方向一致，
其相对竖直的倾角为
\[
\boxed{\delta=\arctan(a_0/g)
\simeq5.82^\circ\approx6^\circ},
\]
从竖直方向朝\emph{支点加速度的反方向}倾斜。
\excheck 求进动速率应用有效重力场的\emph{模}，
求进动轴方向则保留其矢量方向。

\exproblem{286}{46}
质量 $m=5.0\,\mathrm{kg}$、
半径 $R=6.0\,\mathrm{cm}$ 的均匀实心球，
绕穿过球心的水平轴以
$\omega=1250\,\mathrm{rad/s}$ 自转。
其水平轴由两个间距 $\ell=15\,\mathrm{cm}$
的轴承支撑。整个支架绕竖直轴以
$\Omega=5.0\,\mathrm{rad/s}$ 匀速转动。
求轴承受到的陀螺力大小及方向关系。
\exsolution
均匀实心球自转惯量为
$I=\frac25mR^2$，主自转角动量大小 $L=I\omega$。
自转轴水平，其角动量方向随支架绕竖直轴以
角速度 $\Omega$ 转动，
因此角动量方向变化率与其垂直，大小
\[
N_{\rm gyro}=\Omega L=I\omega\Omega
=\frac25mR^2\omega\Omega.
\]
轴承相距 $\ell$，
需要一对大小相等、方向相反的横向支承力
形成这一力偶；每端的陀螺力大小为
\[
\boxed{F=\frac{N_{\rm gyro}}{\ell}
=\frac{2mR^2\omega\Omega}{5\ell}
\simeq3.00\times10^2\,\mathrm N
=0.30\,\mathrm{kN}}.
\]
\excheck 两端轴承反力方向相反，
形成一个力偶，具体上、下方向取决于
自转及支架转动的方向，可由
$\bm N=\bm\Omega\times\bm L$ 判定。

\exproblem{287}{46}
陀螺圆盘质量为 $m=15\,\mathrm{kg}$、
半径 $r=5.0\,\mathrm{cm}$，
自转角速度为 $\omega=330\,\mathrm{rad/s}$。
轴承间距为 $\ell=15\,\mathrm{cm}$。
自转轴被迫绕一水平轴作小幅周期摆动，
摆动周期 $T=1.0\,\mathrm s$、
角振幅 $\varphi_m=20^\circ$。
求轴承所受陀螺力的最大值。
\exsolution
圆盘轴向转动惯量
$I=\frac12mr^2$，
自转主角动量 $L=I\omega$。
设自转轴的摆动角
$\varphi(t)=\varphi_m\sin(2\pi t/T)$，
那么轴线方向改变的瞬时角速度
$|\dot\varphi|$ 最大值为 $2\pi\varphi_m/T$。
维持这种方向变化的陀螺力矩最大值
\[
N_{\max}=L\frac{2\pi\varphi_m}{T}
=\frac{\pi mr^2\omega\varphi_m}{T}.
\]
由轴承距 $\ell$，每端最大附加力
\[
\boxed{F_{\max}
=\frac{\pi mr^2\omega\varphi_m}{\ell T}
\simeq90\,\mathrm N=0.090\,\mathrm{kN}},
\]
其中 $\varphi_m=20^\circ$ 须换算为
$0.349\,\mathrm{rad}$。
\excheck 本题的最大陀螺力发生在角位移过零且
角速度模最大的时刻，而不是角位移达到 $\pm\varphi_m$ 时。

\exproblem{288}{46}
轮船以 $v=36\,\mathrm{km/h}$ 沿水平圆弧航行，
轨道半径为 $R=200\,\mathrm m$。
船上沿船长方向布置的飞轮轴
转动惯量为 $I=3.8\times10^3\,\mathrm{kg\,m^2}$，
飞轮转速 $n=300$ 转/分。
求飞轮轴承承受的陀螺力矩大小。
\exsolution
飞轮自转角速度为
$\omega=2\pi n/60$（弧度每秒）。
船体沿水平圆弧改变航向的角速度为
$\Omega=v/R$，其中
$v=36/3.6=10.0\,\mathrm{m/s}$。
飞轮主自转角动量
$L=I\omega$ 沿船长方向，
与航向改变的竖直角速度 $\bm\Omega$ 垂直。
轴承陀螺力矩大小
\[
\boxed{N=I\omega\Omega
=\frac{2\pi nIv}{60R}
=\frac{2\pi(300)(3.8\times10^3)(10)}{60(200)}
\simeq5.97\times10^3\,\mathrm{N\,m}
\approx6.0\,\mathrm{kN\,m}}.
\]
\excheck 力矩的方向取决于转弯方向和飞轮的自转方向，
可用 $\bm N=\bm\Omega\times\bm L$ 确定。

\exproblem{289}{46}
一台机车沿半径 $R=250\,\mathrm m$ 的水平曲线
以速率 $v=50\,\mathrm{km/h}$ 行驶。
汽轮机转轴平行于车轮轴，其自转方向与车轮相同。
汽轮机转子的轴向惯量
$I=240\,\mathrm{kg\,m^2}$，
转速 $n=1500$ 转/分。两轨轨距
$\ell=1.5\,\mathrm m$。
求陀螺效应使铁轨垂直载荷产生的附加变化量。
\exsolution
机车绕竖直方向转弯角速度
$\Omega=v/R$；
汽轮机转子角速度
$\omega=2\pi n/60$。
转轴水平且与航向改变的竖直轴垂直，
故陀螺力矩大小
$N=I\omega\Omega$。
这一力矩由左右轨道上大小相等、方向相反
的\emph{附加竖直反力}形成，
以轨距 $\ell$ 为力偶臂：
\[
\boxed{\Delta F=\frac{N}{\ell}
=\frac{2\pi nIv}{60R\ell}}.
\]
代入
$v=50/3.6\simeq13.89\,\mathrm{m/s}$：
\[
\boxed{\Delta F\simeq1.40\times10^3\,\mathrm N
=1.4\,\mathrm{kN}}.
\]
\excheck 由于汽轮机自转方向与车轮相同，
按原书图示运动方向，\emph{外轨}载荷增加
$\Delta F$、内轨载荷减少同量；
不是两轨总载荷同时增加。
% === END 05-rigid-body-281-289.tex ===
